
\section{Results and Analysis}
\label{sec:results}

All $F_{0.5}$ and Dice scores are computed at binarization threshold $\theta = 0.4$.

\subsection{Benchmark Comparison}

Table~\ref{tab:benchmark} compares our models against the Grand Prize (GP) Winner on the held-out test segment (20230504223647), never seen during training. The GP Winner uses 26 input layers (17-42); our models use 5 (30--34).

\begin{table*}[t]
\centering
\caption{Benchmark comparison on held-out test segment (20230504223647). All models evaluated at threshold $\theta = 0.4$. Best values in bold. The GP Winner uses 26 input layers (17-42) while our models use 5 layers (30-34).}
\label{tab:benchmark}
\begin{tabular}{llccccccc}
\hline
\textbf{Model} & \textbf{Epochs} & \textbf{$F_{0.5}$} & \textbf{Dice} & \textbf{Precision} & \textbf{Recall} & \textbf{Mean Int.} & \textbf{Cov.@0.4} & \textbf{Layers} \\
\hline
GP Winner TimeSformer  & 30 & 0.379 & 0.332 & 0.418 & 0.275 & 0.298 & 0.174 & 26 \\
\hline
Our UNet2D             & 30  & 0.171 & 0.131 & 0.214 & 0.094 & 0.223 & 0.116 & 5 \\
\textbf{Our UNet2D}    & \textbf{80}  & \textbf{0.474} & \textbf{0.526} & 0.444 & \textbf{0.645} & 0.366 & 0.384 & 5 \\
Our UNet2D             & 100 & 0.341 & 0.326 & 0.352 & 0.303 & 0.270 & 0.228 & 5 \\
\hline
Our TimeSformer        & 30  & 0.346 & 0.230 & \textbf{0.523} & 0.147 & 0.231 & 0.075 & 5 \\
Our TimeSformer        & 80  & 0.403 & 0.311 & 0.504 & 0.225 & 0.256 & 0.118 & 5 \\
Our TimeSformer        & 100 & 0.290 & 0.190 & 0.447 & 0.121 & 0.216 & 0.072 & 5 \\
\hline
\end{tabular}
\end{table*}

The 80-epoch UNet2D achieves $F_{0.5} = 0.474$, outperforming the GP Winner ($F_{0.5} = 0.379$) by $+25\%$ with only 5 input layers $5.2\times$ reduction in input depth. The 80-epoch TimeSformer also surpasses the GP Winner ($F_{0.5} = 0.403$), confirming that the gain is not architecture-specific.

These results require contextualization. The GP Winner was trained on $\sim$20 segments with 15 rounds of expert label refinement~\cite{nader2024gp}, whereas our models train on only 4 segments using pseudo-labels generated by the GP model itself~\cite{lee2013pseudo}. This constitutes a knowledge distillation setup~\cite{hinton2015distilling}: the student (our model) is trained on teacher-generated targets without the teacher's training complexity. The superior recall of our best model (0.645 vs.\ 0.275) reflects focused depth selection-by restricting input to 5 surface-proximal layers, the model avoids the signal dilution inherent in deeper, uninformative layers that lower the GP model's recall. Figure~\ref{fig:benchmark_comparison} shows side-by-side predictions on the test segment.

\begin{figure}[h]
\centering
\includegraphics[width=0.95\textwidth]{benchmark_comparison}
\caption{Side-by-side ink detection on the held-out test segment.}
\label{fig:benchmark_comparison}
\end{figure}

\subsection{Training Duration and Overfitting}
\label{sec:epoch_analysis}

Both architectures exhibit a non-monotonic response to training duration: performance improves from 30 to 80 epochs, then degrades by 100 epochs (Table~\ref{tab:epoch_budget}). Overfitting on a small 4-segment pseudo-label corpus is the primary cause: once the model exhausts the representational diversity of the training tiles, additional epochs shift it toward memorization rather than generalization.

\begin{figure}[h]
\centering
\includegraphics[width=\columnwidth]{training_curves}
\caption{Training loss and validation $F_{0.5}$ for UNet2D and TimeSformer over 100 epochs.}
\label{fig:training_curves}
\end{figure}

\begin{table}[h]
\centering
\caption{Training duration vs.~held-out test performance. Bold = best per architecture.}
\label{tab:epoch_budget}
\begin{tabular}{llcc}
\hline
\textbf{Model} & \textbf{Epochs} & \textbf{Val $F_{0.5}$} & \textbf{Test $F_{0.5}$} \\
\hline
UNet2D & 30  & 0.145 & 0.171 \\
UNet2D & 80  & 0.126 & \textbf{0.474} \\
UNet2D & 100 & 0.144 & 0.341 \\
\hline
TimeSformer & 30  & 0.108 & 0.346 \\
TimeSformer & 80  & 0.127 & \textbf{0.403} \\
TimeSformer & 100 & 0.134 & 0.290 \\
\hline
\end{tabular}
\end{table}

A striking \emph{validation-test divergence} is evident: the 80-epoch UNet2D achieves validation $F_{0.5} = 0.126$ yet test $F_{0.5} = 0.474$, while the 100-epoch model scores \emph{higher} on validation (0.144) but \emph{lower} on test (0.341). The single validation segment is simply not representative of the broader test distribution, so checkpoint selection based on validation alone is unreliable. \textbf{80 epochs is the empirically optimal training budget for this corpus.} All systematic experiments (A--E) use the 100-epoch models as a consistent baseline.

\subsection{Experiment A: Layer-Count Ablation}

The layer count shows a non-monotonic effect: $N = 7$ achieves the best validation $F_{0.5}$ (0.150), followed by $N = 3$ (0.146), $N = 5$ (0.133), and $N = 9$ (0.130). Including layers outside the ink-bearing zone (30--35) introduces uninformative channels that dilute the learned features, consistent with findings on multi-channel medical volume inputs~\cite{roth2014new}. The 100-epoch ImageNet-pre-trained baseline ($N = 5$, $F_{0.5} = 0.144$) outperforms the equivalent ablation model without pre-training (0.133), confirming that transfer learning retains value even for grayscale CT inputs. On balance, $N = 5$ layers represents a practical sweet spot between coverage and signal quality, and is retained as the default configuration.

\subsection{Experiment B: Depth Window Position Sweep}

Table~\ref{tab:exp_b} evaluates five starting positions with $N = 5$ fixed on the 100-epoch UNet2D; model weights are frozen and only the input window shifts.

\begin{table}[h]
\centering
\caption{Experiment B: Position sweep ($N = 5$, metrics averaged over test/val/train segments). Bold = optimal window.}
\label{tab:exp_b}
\begin{tabular}{ccccccc}
\hline
Start $s$ & Layers & $F_{0.5}$ & Dice & Prec. & Recall & Cov.@0.4 \\
\hline
28 & 28--32 & 0.281 & 0.340 & 0.253 & 0.594 & 0.330 \\
29 & 29--33 & 0.281 & 0.329 & 0.259 & 0.535 & 0.264 \\
\textbf{30} & \textbf{30--34} & \textbf{0.293} & \textbf{0.331} & \textbf{0.277} & 0.482 & 0.204 \\
31 & 31--35 & 0.265 & 0.274 & 0.267 & 0.330 & 0.146 \\
32 & 32--36 & 0.274 & 0.245 & 0.311 & 0.239 & 0.103 \\
\hline
\end{tabular}
\end{table}

Window $s = 30$ (layers 30-34) achieves the highest $F_{0.5}$ of 0.293, validating the default training configuration. Depth has a clear precision-recall trade-off: shallower windows over-capture pre-surface tissue, giving high recall but poor precision, while deeper windows miss the ink layer, collapsing recall. Coverage at $\theta = 0.4$ declines monotonically with depth (0.330 $\to$ 0.103), with the sharpest drop between $s = 30$ and $s = 31$ ($-28\%$). This matches the physical expectation: carbon ink applied to the papyrus surface penetrates only a few micrometres~\cite{mocella2015revealing}, so the detectable signal decays rapidly below layer 30.

In contrast, the GP Winner's temporal attention implicitly learns this weighting from 26 input layers. Our approach hard-codes the empirically optimal window, trading flexibility for simplicity.

\subsection{Experiment C: Leave-One-Out Layer Sensitivity}

In reference with Table~\ref{tab:exp_c} each of the five input layers is individually zeroed and the mean absolute change in the prediction map is recorded (100-epoch UNet2D, weights fixed).

\begin{table}[h]
\centering
\caption{Experiment C: Leave-one-out sensitivity (mean $|\Delta\text{Pred}|$ per segment when each layer is masked). Bold = most critical layer.}
\label{tab:exp_c}
\begin{tabular}{ccccc}
\hline
\textbf{Layer} & \textbf{Test} & \textbf{Val} & \textbf{Train} & \textbf{Mean} \\
\hline
30 & 0.160 & 0.130 & 0.098 & 0.129 \\
31 & 0.168 & 0.157 & 0.145 & 0.156 \\
\textbf{32} & \textbf{0.166} & \textbf{0.157} & \textbf{0.153} & \textbf{0.159} \\
33 & 0.143 & 0.141 & 0.104 & 0.129 \\
34 & 0.161 & 0.159 & 0.127 & 0.149 \\
\hline
\end{tabular}
\end{table}

Layers 32 and 31 are most informative (mean sensitivity 0.159 and 0.156), aligning with their position at the papyrus surface where CT density contrast from ink is highest. Crucially, all five layers contribute (range 0.129-0.159) and no layer is dispensable 
masking any single layer causes maximum per-pixel changes of 0.77-0.97, confirming multi-layer corroboration is essential. Outer layers (30 and 34) gain relative importance on unseen segments (layer 30 rises from 0.098 on train to 0.160 on test), suggesting they encode background context that aids generalization.

\subsection{Experiment D: Cross-Segment Generalization}

Table~\ref{tab:exp_d} applies the fixed $N = 5$, $s = 30$ configuration to the 8 ground-truth segments of Scroll~1 using the 100-epoch UNet2D.

\begin{table}[h]
\centering
\caption{Experiment D: Cross-segment results ($N = 5$, $s = 30$) on all 8 ground-truth segments.}
\label{tab:exp_d}
\begin{tabular}{lcccl}
\hline
\textbf{Segment} & \textbf{$F_{0.5}$} & \textbf{Cov.@0.4} & \textbf{Mean Int.} & \textbf{Role} \\
\hline
20230504225948 & 0.393 & 0.438 & 0.431 & Train \\
20230504231922 & 0.394 & 0.156 & 0.258 & Train \\
20230504223647 & 0.338 & 0.225 & 0.268 & Test \\
20230504171956 & 0.308 & 0.075 & 0.215 & Train \\
20230504125349 & 0.283 & 0.266 & 0.322 & Train \\
20230504093154 & 0.147 & 0.231 & 0.301 & Validation \\
20230503225234 & 0.000 & 0.151 & 0.186 & \textbf{Failure} \\
20230504094316 & 0.000 & 0.001 & 0.160 & \textbf{Failure} \\
\hline
\end{tabular}
\end{table}

$F_{0.5}$ spans 0.000-0.394 with a clear bimodal structure: six segments succeed (0.147-0.394) and two fail completely, yielding a 25\% failure rate. The held-out test segment (0.338) falls within the training range (0.283-0.394), indicating that the model has learned transferable features rather than memorizing training tiles.

\subsection{Experiment E: Failure Analysis}

The two complete failures from Experiment~D were diagnosed by examining prediction coverage and spatial overlap with ground truth.

\textbf{Segment 20230504094316} coverage $= 0.001$ ($\approx 0.1\%$). The model detects no ink. The cause is a \emph{depth offset}: the true papyrus surface in this segment lies outside the fixed layer 30-34 window, placing the ink signal beyond the model's input range. Per-segment intensity profiling to shift the window to the actual surface peak would directly resolve this.

\textbf{Segment 20230503225234} coverage $= 0.151$ but $F_{0.5} = 0.000$. The model produces spatially plausible predictions, but with zero overlap against ground truth. This indicates \emph{spatial misalignment}, most likely an unregistered Z-axis reversal or errors in the pseudo-label for this segment.

The key insight is that both failures are \emph{data-driven}, not model-driven: the underlying features that make ink detectable are present, but incorrect depth placement or label misalignment prevents correct localization. A static global Z-window is sufficient for most segments; robust scroll-scale deployment requires per-segment depth profiling to guarantee the ink signal falls within the selected window.

\vspace{2\baselineskip}

\newpage